\chapter{Implementacja}\label{ch:implementacja}

% ═══════════════════════════════════════════════════════════════════════════════
% Cel rozdziału: 15–25 stron. Szczegóły implementacyjne.
% Struktura odzwierciedla monorepo: dataset_pipeline, on_premise, cloud,
% gateway, frontend.  Każda sekcja: decyzje projektowe + kluczowe pliki
% + fragmenty kodu (listingi) + problemy napotkane i rozwiązania.
% ═══════════════════════════════════════════════════════════════════════════════

\section{Organizacja monorepo}\label{sec:monorepo}

% TREŚĆ DO NAPISANIA:
%
% 1. Struktura katalogów (drzewko):
%    MedVoicePL/
%    ├── dataset_pipeline/    (workspace member, src-layout)
%    ├── application/
%    │   ├── on_premise/      (workspace member)
%    │   ├── cloud/           (workspace member)
%    │   ├── gateway/         (workspace member)
%    │   ├── frontend/        (Node.js, poza uv workspace)
%    │   └── shared/          (workspace member, współdzielone schematy)
%    ├── research_and_testing/ (workspace member)
%    ├── pyproject.toml       (root — workspace definition, linting, testing)
%    └── docker-compose.yml
%
% 2. Zarządzanie zależnościami:
%    - uv workspace: [tool.uv.workspace] members = [...]
%    - src-layout: <member>/src/<pkg>/ — unika kolizji top-level `src`.
%    - Importy: `from dataset_pipeline.pipeline import run_audio`.
%    - Problem rozwiązany: kolizja pakietów (każdy member instalował się
%      jako `src`; restrukturyzacja na `<member>/src/<pkg>/`).
%
% 3. Współdzielone schematy (shared/):
%    - Pydantic v2 modele jako kontrakty między modułami.
%    - AnonymizedPayload, SOAPNote, DrugSuggestion.
%    - Walidacja na granicy każdego modułu.
%
% 4. Konfiguracja jakości kodu:
%    - ruff: select E/W/F/I/N/UP/B/SIM/TCH/RUF, target-version py311.
%    - mypy: python_version 3.11, ignore_missing_imports (relaxed for MVP).
%    - pytest: testpaths per moduł, addopts -v --tb=short.
%
% LISTING: fragment pyproject.toml — workspace definition.

\section{Pipeline datasetu syntetycznego}\label{sec:dataset-pipeline}

% TREŚĆ DO NAPISANIA — INTRO:
%
% Cel pipeline'u: wygenerować syntetyczny dataset dialogów medycznych PL
% jako ground truth dla ewaluacji STT, SLM i anonimizacji.
% Trzy fazy: 1a (scenariusze JSON), 1b (synteza mowy), 1b.1 (augmentacje).

\subsection{Generacja scenariuszy (Phase 1a)}\label{sec:scenariusze}

% TREŚĆ DO NAPISANIA:
%
% 1. Provider LLM: Google Gemini przez Vertex AI.
%    - SDK: google-genai (vertexai=True).
%    - Auth: Application Default Credentials (ADC).
%    - Structured output: response_schema=ScenarioModel → JSON walidowany
%      przez Pydantic natychmiast (nie parsowanie tekstu).
%
% 2. Kontrakt danych — ScenarioModel (schema.py):
%    - ScenarioModel: id, specialization, patient_sex, chief_complaint,
%      dialogue (list[DialogueTurn]), ground_truth_entities.
%    - DialogueTurn: speaker (SpeakerRole), text, turn_index.
%    - Specialization enum: CARDIOLOGY, ENDOCRINOLOGY, NEUROLOGY,
%      PSYCHIATRY, ORTHOPEDICS.
%
% 3. Prompt engineering (prompts.py):
%    - SYSTEM_PROMPT_PL: kontekst lekarza POZ, język polski, specyfika
%      dokumentacji medycznej, wymogi RODO (fikcyjne PII).
%    - FEW_SHOT_EXAMPLES: 5 ręcznie napisanych scenariuszy (1 per spec.).
%    - render_user_prompt: dynamiczne parametry (specjalizacja, scenariusz nr).
%
% 4. Generator (generator.py):
%    - ScenarioGenerator: batch + ThreadPoolExecutor (4 wątki).
%    - Retry: tenacity (exp backoff, 4 próby) + correction hint retry
%      (2 próby przy błędach walidacji Pydantic).
%    - Output: scenarios/*.json + failures.json.
%
% 5. Skala: 8 scenariuszy (1-2 per specjalizacja × 5).
%
% LISTINGI: fragment ScenarioModel, fragment SYSTEM_PROMPT_PL.

\subsection{Synteza mowy (Phase 1b)}\label{sec:tts}

% TREŚĆ DO NAPISANIA:
%
% 1. Model TTS: XTTSv2 przez community fork coqui-tts (>=0.27).
%    - Architektura: VITS-like, voice cloning, wielojęzyczność.
%    - Model ~2 GB, cache pod TTS_MODEL_CACHE.
%    - Autodetect: MPS (Apple Silicon) / CUDA / CPU.
%
% 2. Próbki referencyjne głosów (voice_manager.py):
%    - Źródło: Common Voice Scripted Speech 25.0 — Polish (Mozilla Data
%      Collective, Bearer auth).
%    - 4 profile: doctor_male, doctor_female, patient_male, patient_female.
%    - Filtr: długość zdania 130–260 znaków (≈10–20s) → stabilna prozodia.
%    - Relaxation ladder: 90, potem 60 znaków gdy filtry za ciasne.
%    - Skrypt fetch_reference_voices.py (idempotentny).
%
% 3. Stabilność syntazy (xtts_synthesizer.py):
%    - temperature=0.65 (poniżej domyślnego 0.75).
%    - sound_norm_refs=True (auto-normalizacja referencji).
%    - split_sentences=True (zdanie-po-zdaniu → lepsza prozodia PL).
%    - Czyszczenie tekstu PL: usunięcie artefaktów prompt, normalizacja.
%
% 4. Determinizm (kluczowy dla reprodukowalności):
%    - Seed miksera: sha256(scenario.id)[:8].
%    - Seed per-turn: sha256(scenario.id)[:4] XOR (turn_index * 0x9E3779B1).
%    - Trzy RNG: torch.manual_seed + np.random.seed + random.seed
%      (coqui-tts sampluje z trzech niezależnych RNG).
%    - Weryfikacja: byte-identyczny dialogue.wav między reruns na tym samym HW.
%
% 5. Postprodukcja (dialog_mixer.py):
%    - Equal-power crossfade między turami.
%    - Randomized pause (seeded) między turami.
%    - Resample: per-turn 22050 Hz → mix 16 kHz.
%    - Normalize (peak-normalize).
%
% 6. Output per scenariusz:
%    - turn_NN_<speaker>.wav (22050 Hz) — ground truth per tura.
%    - dialogue.wav (16 kHz) — clean mix, input dla STT.
%    - bundle.json — metadane (SHA, duracje, parametry).
%
% LISTINGI: fragment XTTSSynthesizer.synthesize(), MixerConfig.
% RYSUNEK: diagram pipeline'u Phase 1b (flow chart).

\subsection{Augmentacja i szum (Phase 1b.1)}\label{sec:augmentacja}

% TREŚĆ DO NAPISANIA:
%
% 1. Cel: realistyczna degradacja „clinic ambience" obok czystego dialogu.
%
% 2. Per-turn augmentacja (augmentation.py, pedalboard):
%    - Pitch shift: ±1 semitone.
%    - Time stretch: 0.95–1.05.
%    - Aplikowane PRZED miksowaniem → naturalna wariancja.
%
% 3. Dialog-level reverb (augmentation.py):
%    - Rozmiary pokoju: small=0.15, medium=0.35, large=0.55.
%    - Wagi: 0.4 / 0.4 / 0.2.
%    - Aplikowane RAZ na zmiksowanym dialogu.
%
% 4. Noise injection (noise_injector.py):
%    - Biblioteka szumów: Freesound.org (CC0/CC-BY), 7 kategorii
%      (keyboard, HVAC, paper, chatter, footsteps, door, hospital).
%    - SNR stratified buckets: 20–25 dB (40\%), 15–20 dB (40\%), 10–15 dB (20\%).
%    - RMS-based gain matching.
%    - Seed: sha256(scenario.id + "|noise")[:8] — oddzielony od clean.
%
% 5. Kolejność efektów: mix → reverb → resample 16kHz → noise → normalize.
%
% 6. Output: dialogue.wav (clean, nietknięty) + dialogue\_noisy.wav.
%
% LISTING: fragment NoiseConfig, SnrBucket.

\section{Moduł on-premise}\label{sec:impl-on-premise}

\subsection{Transkrypcja audio}\label{sec:stt-impl}

% TREŚĆ DO NAPISANIA:
%
% 1. faster-whisper: CTranslate2 backend, model large-v3.
%    - Kwantyzacja: int8 (CPU), float16 (CUDA). MPS nie wspierane
%      przez ctranslate2 → CPU fallback na Apple Silicon.
%    - Alternatywa Apple Silicon: mlx-whisper (Metal GPU, ~5–10× szybciej).
%
% 2. Diaryzacja: pyannote.audio (VAD + embedding + clustering).
%    - Dialog 2-osobowy: lekarz + pacjent.
%    - Mapowanie: embedding cluster → rola (speaker_0, speaker_1).
%
% 3. Output: lista segmentów z timestampem, etykietą mówcy, tekstem.
%
% 4. Wydajność (Apple Silicon M5, large-v3):
%    - int8 CPU: RTF ~0.6 (3 min dialog = ~5 min).
%    - MLX Metal: RTF ~0.1 (3 min dialog = ~1–2 min).
%
% LISTING: fragment konfiguracji faster-whisper.

\subsection{Strukturyzacja SOAP}\label{sec:slm-impl}

% TREŚĆ DO NAPISANIA:
%
% 1. Wybór modelu SLM: [do określenia w Phase 2 — Phi-3/Llama 3.x/Bielik].
% 2. Prompt engineering: system prompt + transkrypcja → JSON SOAP.
% 3. Kwantyzacja: GGUF 4-bit, ładowanie przez transformers/llama.cpp.
% 4. Output: SOAPNote Pydantic model (S/O/A/P fields).
% 5. Ewaluacja jakości: ekspert kliniczny, metryki spójności.

\subsection{Anonimizacja PII}\label{sec:presidio-impl}

% TREŚĆ DO NAPISANIA:
%
% 1. Microsoft Presidio — Analyzer + Anonymizer:
%    - Wbudowane recognizery: spaCy NER (pl_core_news_lg).
%    - Custom recognizery PL:
%      a) PeselRecognizer: regex \d{11} + walidacja cyfry kontrolnej.
%      b) NipRecognizer: regex \d{10} + walidacja.
%      c) PolishPhoneRecognizer: +48 / 0XX / XXX-XXX-XXX.
%      d) PolishPostalCodeRecognizer: XX-XXX.
%      e) PolishNameRecognizer: deny-list polskich imion.
%    - Strategia: replace (np. [PESEL], [IMIĘ], [TELEFON]).
%
% 2. trust_gate.py:
%    - Walidacja payload wychodzącego z on-premise.
%    - Regex scan na znane wzorce PII (PESEL, NIP, telefon).
%    - Blokada + alert jeśli wykryty potencjalny wyciek.
%
% 3. audit_logger.py:
%    - structlog z timestampem, session_id, typem encji, akcją.
%    - Logi lokalne, nigdy transmitowane zewnętrznie.
%
% LISTING: fragment custom PeselRecognizer, trust_gate.py walidacja.

\section{Moduł chmurowy}\label{sec:impl-cloud}

\subsection{Analiza kliniczna LLM}\label{sec:llm-impl}

% TREŚĆ DO NAPISANIA:
%
% 1. Provider: OpenAI / Anthropic / Google (konfigurowalne).
% 2. Structured output (JSON mode / function calling).
% 3. Prompty kliniczne: kontekst specjalizacji, guardrails (ograniczenie
%    halucynacji), zakaz diagnozowania — tylko sugestie.
% 4. Input: AnonymizedPayload. Output: ClinicalAnalysis Pydantic model.

\subsection{Retrieval-Augmented Generation}\label{sec:rag-impl}

% TREŚĆ DO NAPISANIA:
%
% 1. ChromaDB w trybie embedded (MVP, bez zewnętrznego serwera).
% 2. Embeddingi: multilingual model (mE5/BGE-M3).
% 3. Baza wiedzy: leki refundowane NFZ, ChPL (charakterystyki produktów
%    leczniczych), wytyczne towarzystw naukowych.
% 4. Chunkowanie: recursive text splitter, 512 tokenów, overlap 64.
% 5. Retrieval: top-k dense + opcjonalnie BM25 sparse → reranking.
% 6. Output: lista pasujących dokumentów z score i źródłem.

\subsection{Dopasowanie refundacji NFZ}\label{sec:nfz-impl}

% TREŚĆ DO NAPISANIA:
%
% 1. Import bazy refundacji NFZ (CSV/Excel publikowany przez MZ/NFZ).
% 2. Pola: nazwa międzynarodowa, nazwa handlowa, dawka, postać, cena,
%    kategoria refundacji (R, Rp, S, Rpz), odpłatność (30\%, 50\%, ryczałt).
% 3. Wyszukiwanie rozmyte: Levenshtein / fuzzy matching nazw leków.
% 4. Dopasowanie: leki sugerowane przez LLM → najlepsze refundowane
%    odpowiedniki z bazy NFZ.
% 5. Output: DrugSuggestion Pydantic model (lek, refundacja, cena, źródło).

\section{Gateway i frontend}\label{sec:impl-gateway-frontend}

% TREŚĆ DO NAPISANIA:
%
% 1. Gateway (FastAPI):
%    - Endpointy REST (z api_reference.md).
%    - WebSocket dla opcjonalnego streamingu statusu.
%    - Routing: on-premise → anonimizacja → cloud → wynik.
%    - Health checks: /health per serwis.
%    - Middleware: CORS, rate limiting, request logging (structlog).
%
% 2. Frontend (Next.js + TypeScript):
%    - Nagrywanie audio: MediaRecorder API → WAV blob → upload.
%    - Dashboard: transkrypcja, SOAP, sugestie leków, status.
%    - Server Actions: formularz akceptacji/odrzucenia sugestii.
%    - Responsywność: desktop (monitor w gabinecie) + tablet.

\section{Fragment implementacji --- przykład}\label{sec:kod-przyklad}

Poniższy listing prezentuje kontrakt Pydantic używany jako granica
przesyłania danych między modułem lokalnym a~chmurą.

\begin{lstlisting}[language=Python,caption={Kontrakt AnonymizedPayload (uproszczony)},label={lst:payload}]
from pydantic import BaseModel, Field

class AnonymizedPayload(BaseModel):
    """Jedyna struktura przekraczająca granicę zaufania."""
    session_id: str = Field(..., description="Identyfikator sesji")
    soap_note: str = Field(..., description="Notatka w formacie SOAP")
    pii_detected: list[str] = Field(default_factory=list)
    audit_hash: str = Field(..., description="Hash SHA-256 oryginału")
\end{lstlisting}

% DODATKOWE LISTINGI DO DODANIA:
%   - ScenarioModel (schema.py) — kontrakt datasetu.
%   - Fragment GeminiClient (llm_client.py) — structured output.
%   - Fragment XTTSSynthesizer (xtts_synthesizer.py) — seed + synteza.
%   - Fragment trust_gate.py — walidacja PII.
%   - Fragment NoiseConfig — SNR buckets.
